\chapter{آرشیو‌سازی و پشتیبان‌گیری}

یکی از مسئولیت‌های بنیادین هر مدیر سیستم (\en{SysAdmin})، تضمین امنیت و پایداری داده‌های موجود در سیستم است. کلیدی‌ترین راهبرد برای تحقق این هدف، اجرای منظم عملیات پشتیبان‌گیری (\en{Backup}) از فایل‌های حیاتی است. با این حال، حتی برای کاربران عادی نیز تهیه نسخه‌های پشتیبان، جابه‌جایی مجموعه‌های حجیم از داده‌ها میان مسیرهای مختلف و انتقال فایل‌ها بین دستگاه‌های گوناگون، مهارتی بسیار کاربردی محسوب می‌شود.

در این فصل، به بررسی ابزارهای استاندارد و پرکاربردی خواهیم پرداخت که جهت مدیریت مجموعه‌ای از فایل‌ها توسعه یافته‌اند. این ابزارها در سه دسته‌ی کلی طبقه‌بندی می‌شوند؛ ابتدا برنامه‌های فشرده‌سازی داده‌ها:

\begin{itemize}
  \item \cmd{gzip}: ابزاری جهت فشرده‌سازی (\en{Compress}) یا گشودن (\en{Expand}) فایل‌ها.
  \item \cmd{bzip2}: فشرده‌ساز فایل مبتنی بر الگوریتم مرتب‌سازی بلوکی (\en{Block-sorting}).
\end{itemize}

سپس ابزارهای آرشیو‌سازی (\en{Archiving}) را بررسی می‌کنیم:

\begin{itemize}
  \item \cmd{tar}: ابزار کلاسیک آرشیو نوار مغناطیسی (\en{Tape Archive}) که برای تجمیع چندین فایل در یک فایل واحد به کار می‌رود.
  \item \cmd{zip}: ابزاری دو منظوره جهت بسته‌بندی و فشرده‌سازی هم‌زمان فایل‌ها (سازگار با محیط ویندوز).
\end{itemize}

و در نهایت، ابزار همگام‌سازی داده‌ها:

\begin{itemize}
  \item \cmd{rsync}: ابزاری قدرتمند برای همگام‌سازی (\en{Synchronization}) بهینه فایل‌ها و دایرکتوری‌ها، به‌ویژه در ارتباطات راه دور (\en{Remote}).
\end{itemize}

\section{مبانی فشرده‌سازی فایل}

در طول تاریخ رایانش، همواره تلاشی مستمر برای جای‌گذاری بیشترین حجم داده در کمترین فضای ممکن در جریان بوده است؛ خواه این فضا مربوط به حافظه و ذخیره‌سازها باشد و خواه پهنای باند شبکه. بسیاری از سرویس‌های دیجیتالی مدرن که امروزه به آن‌ها متکی هستیم، نظیر شبکه‌های تلفن همراه، تلویزیون‌های با وضوح بالا (\en{HDTV}) و اینترنت پرسرعت، همگی مدیون تکنیک‌های پیشرفته‌ی فشرده‌سازی داده‌ها هستند.

فشرده‌سازی داده (\en{Data Compression}) در واقع فرآیند شناسایی و حذف افزونگی (\en{Redundancy}) از اطلاعات است. برای درک بهتر، یک مثال انتزاعی را در نظر بگیرید: فرض کنید تصویری کاملاً سیاه با ابعاد ۱۰۰ در ۱۰۰ پیکسل داریم. در حالت ذخیره‌سازی خام (\en{Raw}) با فرض اختصاص ۲۴ بیت (۳ بایت) به هر پیکسل، این تصویر ۳۰,۰۰۰ بایت فضا اشغال خواهد کرد:

\begin{latin}
\begin{lstlisting}
100 x 100 x 3 = 30,000
\end{lstlisting}
\end{latin}

تصویری که تنها از یک رنگ تشکیل شده، دارای داده‌های کاملاً افزونه است. به جای ذخیره‌ی بلوکی حاوی ۳۰,۰۰۰ بایت صفر (که معمولاً نشان‌دهنده‌ی رنگ سیاه است)، می‌توانیم به شکلی هوشمندانه فقط توصیف کنیم که «یک بلوک شامل ۱۰,۰۰۰ پیکسل سیاه» داریم. این روش که داده‌ها را به یک عدد (تعداد تکرار) و خود داده (رنگ سیاه) خلاصه می‌کند، «کدگذاری طول اجرا» (\en{Run-Length Encoding}) نامیده می‌شود و از ابتدایی‌ترین تکنیک‌های این حوزه به شمار می‌رود. هدف در تمامی روش‌های پیشرفته‌ی امروزی نیز ثابت است: کاهش حجم از طریق حذف داده‌های تکراری.

الگوریتم‌های فشرده‌سازی (روش‌های ریاضیاتی اجرای این فرآیند) به‌طور کلی به دو دسته‌ی اصلی تقسیم می‌شوند:

\begin{itemize}
  \item \textbf{بدون اتلاف (\en{Lossless}):} در این روش، تمامی داده‌های موجود در فایل اصلی حفظ می‌شوند. این بدان معناست که پس از بازگشت فایل از حالت فشرده (\en{Decompression})، نسخه‌ی بازسازی‌شده دقیقاً و بیت‌به‌بیت با نسخه‌ی اولیه یکسان است.

  \item \textbf{با اتلاف (\en{Lossy}):} در این روش برای دستیابی به نرخ فشرده‌سازی بسیار بالاتر، بخشی از داده‌های غیرضروری حذف می‌شوند. در نتیجه، نسخه‌ی بازسازی‌شده دقیقاً مشابه اصل نیست، بلکه تقریب بسیار نزدیکی از آن است. نمونه‌های بارز این تکنیک شامل فرمت \en{JPEG} (برای تصاویر) و \en{MP3} (برای صوت) هستند.
\end{itemize}

در این مبحث، تمرکز ما منحصراً بر فشرده‌سازی بدون اتلاف خواهد بود؛ چرا که اکثر داده‌های سیستمی و متنی در کامپیوترها، حساسیت بالایی داشته و تحمل از دست رفتن حتی یک بیت از اطلاعات را ندارند.

\section{فشرده‌سازی فایل با ابزار gzip}

برنامه \cmd{gzip} جهت فشرده‌سازی یک یا چند فایل طراحی شده است. این ابزار در حالت پیش‌فرض، فایل اصلی را با نسخه‌ی فشرده‌ی آن جایگزین می‌کند. ابزار مکمل آن، \cmd{gunzip} نام دارد که برای بازگرداندن فایل‌های فشرده (\en{Decompression}) به حالت اولیه و بدون فشرده‌سازی به کار می‌رود. به مثال کاربردی زیر توجه کنید:

\begin{latin}
\begin{lstlisting}
[me@linuxbox ~]$ ls -l /etc > foo.txt
[me@linuxbox ~]$ ls -l foo.*
-rw-r--r-- 1 me   me   15738 2025-10-14 07:15 foo.txt
[me@linuxbox ~]$ gzip foo.txt
[me@linuxbox ~]$ ls -l foo.*
-rw-r--r-- 1 me   me    3230 2025-10-14 07:15 foo.txt.gz
[me@linuxbox ~]$ gunzip foo.txt
[me@linuxbox ~]$ ls -l foo.*
-rw-r--r-- 1 me   me   15738 2025-10-14 07:15 foo.txt
\end{lstlisting}
\end{latin}

در این سناریو، ابتدا یک فایل متنی به نام \file{foo.txt} ایجاد کردیم. با اجرای \cmd{gzip}، فایل اصلی حذف و نسخه‌ای با پسوند \file{.gz} جایگزین شد. مشاهده می‌کنید که حجم فایل فشرده تقریباً به یک‌پنجم کاهش یافته، در حالی که مجوزهای دسترسی (\en{Permissions}) و برچسب زمانی (\en{Timestamp}) فایل اصلی کاملاً حفظ شده است. در نهایت با اجرای \cmd{gunzip}، فایل به وضعیت اصلی خود بازگشت.

ابزار \cmd{gzip} از پارامترهای متنوعی پشتیبانی می‌کند که در جدول زیر تشریح شده‌اند:

\begin{table}[htbp]
\centering
\caption{پارامترهای عملیاتی دستور \cmd{gzip}}
\label{tab:gzip-options}
\begin{tabularx}{\textwidth}{l l X}
\toprule
\textbf{گزینه کوتاه} & \textbf{گزینه بلند} & \textbf{شرح عملکرد فنی} \\
\midrule
\cmd{-c} & \cmd{--stdout} \cmd{--to-stdout} & خروجی فشرده را به خروجی استاندارد (\en{stdout}) ارسال کرده و فایل اصلی را بدون تغییر حفظ می‌کند. اگر چندین فایل ورودی مشخص کنید، \cmd{gzip} هر یک را به‌طور جداگانه فشرده کرده و سپس خروجی فشرده‌شده‌ی همه‌ی آن‌ها را پشت سر هم قرار می‌دهد. در نتیجه، یک جریان داده‌ی واحد حاصل می‌شود که شامل چندین بخش فشرده‌شده‌ی مجزا است. \\
\addlinespace
\cmd{-d} & \cmd{--decompress} \cmd{--uncompress} & فایل را از حالت فشرده خارج (\en{decompress}) می‌کند. عملکردی معادل \cmd{gunzip} دارد. \\
\addlinespace
\cmd{-f} & \cmd{--force} & فشرده‌سازی را حتی در صورت وجود نسخه‌ی فشرده‌شده از فایل اصلی، اجباراً انجام می‌دهد. \\
\addlinespace
\cmd{-h} & \cmd{--help} & خلاصه‌ای از نحوه‌ی استفاده و راهنمای گزینه‌ها را نمایش می‌دهد. \\
\addlinespace
\cmd{-l} & \cmd{--list} & اطلاعاتی از جمله حجم فایل فشرده‌شده، حجم فایل اصلی (پیش از فشرده‌سازی)، نسبت فشرده‌سازی و نام فایل را فهرست می‌کند. در ترکیب با \cmd{-v}، جزئیات بیشتری مانند روش فشرده‌سازی، \en{CRC} و زمان نیز نمایش داده می‌شود. \\
\addlinespace
\cmd{-r} & \cmd{--recursive} & پیمایش بازگشتی دایرکتوری‌ها و فشرده‌سازی تمامی فایل‌های درون آن‌ها. اگر یکی از فایل‌های ورودی یک دایرکتوری باشد، \cmd{gzip} به آن وارد شده و تمام فایل‌های موجود در آن را فشرده می‌کند. \\
\addlinespace
\cmd{-t} & \cmd{--test} & یکپارچگی فایل فشرده را بررسی کرده و سلامت ساختار آن را اعتبارسنجی می‌کند. \\
\addlinespace
\cmd{-v} & \cmd{--verbose} & نام هر فایل و درصد کاهش حجم را نمایش می‌دهد. در ترکیب با \cmd{-l}، جزئیات بیشتری مانند روش فشرده‌سازی، \en{CRC} و زمان نیز نمایش داده می‌شود. \\
\addlinespace
\cmd{-number} & --- & تعیین سطح فشرده‌سازی از \cmd{1} (کمترین زمان، کمترین تراکم) تا \cmd{9} (بیشترین زمان، بیشترین تراکم). مقادیر \cmd{1} و \cmd{9} معادل \cmd{--fast} و \cmd{--best} هستند. مقدار پیش‌فرض \cmd{6} است. \\
\bottomrule
\end{tabularx}
\end{table}

با بازگشت به مثال پیشین، سناریوی زیر را در نظر بگیرید:

\begin{latin}
\begin{lstlisting}
[me@linuxbox ~]$ gzip foo.txt
[me@linuxbox ~]$ gzip -tv foo.txt.gz
foo.txt.gz:    OK
[me@linuxbox ~]$ gzip -d foo.txt.gz
\end{lstlisting}
\end{latin}

در این نمونه، ابتدا فایل \file{foo.txt} را به نسخه‌ی فشرده‌ی \file{foo.txt.gz} تبدیل کردیم. سپس با ترکیب گزینه‌های \cmd{-t} (تست سلامت) و \cmd{-v} (گزارش جزئیات)، از صحت ساختار فایل فشرده اطمینان حاصل کرده و در نهایت با پارامتر \cmd{-d}، آن را به حالت اولیه بازگرداندیم.

یکی از قابلیت‌های کلیدی \cmd{gzip}، تعامل بهینه با جریان‌های ورودی و خروجی استاندارد (\en{standard I/O streams}) است:

\begin{latin}
\begin{lstlisting}
[me@linuxbox ~]$ ls -l /etc | gzip > foo.txt.gz
\end{lstlisting}
\end{latin}

این دستور خروجی متنی حاصل از فهرست‌بندی دایرکتوری را مستقیماً از طریق پایپ (\en{Pipe}) دریافت کرده و نسخه‌ی فشرده‌ی آن را در قالب یک فایل ذخیره می‌کند.

ابزار \cmd{gunzip} به‌طور پیش‌فرض فرض می‌کند که فایل‌های هدف دارای پسوند \file{.gz} هستند؛ از این رو، تا زمانی که نام فایل با نسخه‌ی بازشده‌ی آن در همان مسیر تداخل نداشته باشد، نیازی به درج پسوند در دستور نیست:

\begin{latin}
\begin{lstlisting}
[me@linuxbox ~]$ gunzip foo.txt
\end{lstlisting}
\end{latin}

چنانچه صرفاً قصد مشاهده‌ی محتوای یک فایل متنی فشرده را بدون خارج کردن آن از حالت فشرده داشته باشید، می‌توانید از ترکیب زیر استفاده کنید:

\begin{latin}
\begin{lstlisting}
[me@linuxbox ~]$ gunzip -c foo.txt | less
\end{lstlisting}
\end{latin}

همچنین ابزاری به نام \cmd{zcat} در بسته‌ی \cmd{gzip} ارائه شده است که عملکردی معادل \cmd{gunzip -c} دارد و به شما اجازه می‌دهد مانند دستور \cmd{cat}، محتوای فایل‌های فشرده را مستقیماً به خروجی استاندارد بفرستید:

\begin{latin}
\begin{lstlisting}
[me@linuxbox ~]$ zcat foo.txt.gz | less
\end{lstlisting}
\end{latin}

\begin{note}
\textbf{نکته:} برای سهولت بیشتر، ابزاری به نام \cmd{zless} نیز در سیستم‌های لینوکسی موجود است که دقیقاً مشابه پایپ‌لاین فوق عمل کرده و محتوای فایل فشرده را مستقیماً در محیط \cmd{less} باز می‌کند.
\end{note}

\section{فشرده‌سازی فایل با ابزار bzip2}

برنامه \cmd{bzip2} که توسط جولیان سوارد (\en{Julian Seward}) ساخته شده، مشابه \cmd{gzip} است اما از الگوریتم فشرده‌سازی متفاوتی استفاده می‌کند که سطوح فشرده‌سازی بالاتری را با هزینه سرعت کمتر فشرده‌سازی ارائه می‌دهد. از نظر عملکرد نیز، به همان شیوه \cmd{gzip} عمل می‌کند. فایلی که با \cmd{bzip2} فشرده شده باشد با پسوند \file{.bz2} مشخص می‌شود.

\begin{latin}
\begin{lstlisting}
[me@linuxbox ~]$ ls -l /etc > foo.txt
[me@linuxbox ~]$ ls -l foo.txt
-rw-r--r-- 1 me   me   15738 2025-10-17 13:51 foo.txt
[me@linuxbox ~]$ bzip2 foo.txt
[me@linuxbox ~]$ ls -l foo.txt.bz2
-rw-r--r-- 1 me   me    2792 2025-10-17 13:51 foo.txt.bz2
[me@linuxbox ~]$ bunzip2 foo.txt.bz2
\end{lstlisting}
\end{latin}

همان‌طور که مشاهده می‌شود، ابزار \cmd{bzip2} از نظر الگوی استفاده شباهت بسیاری به \cmd{gzip} دارد. این برنامه تقریباً از تمامی پارامترهای عملیاتی \cmd{gzip} (به استثنای گزینه \cmd{-r}) پشتیبانی می‌کند. با این حال، باید توجه داشت که پارامتر تعیین سطح فشرده‌سازی (\cmd{-number}) در \cmd{bzip2} از نظر محاسباتی و تخصیص حافظه، معنای متفاوتی نسبت به رقیب خود دارد.

در بسته‌ی نرم‌افزاری \cmd{bzip2}، ابزارهای \cmd{bunzip2} جهت گشودن فایل‌ها و \cmd{bzcat} برای ارسال مستقیم محتوای فشرده به خروجی استاندارد ارائه شده‌اند. همچنین، این مجموعه شامل ابزار ارزشمندی به نام \cmd{bzip2recover} است که وظیفه دارد داده‌های قابل بازیابی را از فایل‌های آسیب‌دیده‌ی \file{.bz2} استخراج کند.

\section{توصیه فنی: پرهیز از فشرده‌سازی مضاعف}

یک اشتباه رایج در میان کاربران، تلاش برای فشرده‌سازی مجدد فایل‌هایی است که پیش‌تر با الگوریتم‌های کارآمد فشرده شده‌اند؛ برای نمونه، فشرده‌سازی یک تصویر \en{JPEG} با دستور:

\begin{latin}
\begin{lstlisting}
$ gzip picture.jpg
\end{lstlisting}
\end{latin}

نه تنها منجر به هدررفت زمان و منابع پردازشی می‌شود، بلکه در اغلب موارد حجم فایل نهایی را نسبت به نسخه‌ی اولیه افزایش می‌دهد.

دلیل این پدیده آن است که تمامی تکنیک‌های فشرده‌سازی دارای مقداری «سربار» (\en{Overhead}) هستند که شامل داده‌های توصیفی الگوریتم برای بازسازی فایل است. زمانی که سعی می‌کنید فایلی را که فاقد داده‌های افزونه (\en{Redundancy}) است فشرده کنید، الگوریتم فضایی برای کاهش حجم پیدا نمی‌کند و در نتیجه، اضافه‌شدن داده‌های مربوط به سربار، منجر به بزرگ‌تر شدن حجم فایل خروجی می‌گردد.

\section{اصول و مفاهیم آرشیو‌سازی فایل}

یکی از عملیات‌های کلیدی در مدیریت فایل که غالباً با فشرده‌سازی پیوند می‌خورد، آرشیو‌سازی (\en{Archiving}) است. آرشیو‌سازی به فرآیند تجمیع چندین فایل مجزا و بسته‌بندی آن‌ها در قالب یک فایل واحد و بزرگ اطلاق می‌شود.

این فرآیند به‌طور معمول به عنوان رکنی اساسی در عملیات پشتیبان‌گیری (\en{Backup}) سیستم شناخته می‌شود. همچنین، زمانی که نیاز است داده‌های قدیمی و کم‌کاربرد جهت آزادسازی فضا به حافظه‌های ذخیره‌سازی بلندمدت منتقل شوند، از روش آرشیو‌سازی برای یکپارچه‌سازی و سهولت در جابه‌جایی استفاده می‌گردد.

\section{مدیریت آرشیو با ابزار tar}

در اکوسیستم سیستم‌عامل‌های شبه‌یونیکس، برنامه \cmd{tar} ابزار کلاسیک و استاندارد جهت آرشیو‌سازی فایل‌ها محسوب می‌شود. نام این ابزار که سرواژه‌ی \en{Tape Archive} (آرشیو نوار مغناطیسی) است، به ریشه‌های تاریخی آن در ایجاد نسخه‌های پشتیبان روی نوارهای مغناطیسی اشاره دارد. اگرچه \cmd{tar} همچنان برای مقاصد سنتی به کار می‌رود، اما امروزه برای مدیریت آرشیوها بر روی انواع رسانه‌های ذخیره‌سازی مدرن نیز به همان اندازه کارآمد است.

در دنیای لینوکس، غالباً با فایل‌هایی مواجه می‌شوید که دارای پسوند \file{.tar} یا \file{.tgz} هستند؛ پسوند نخست نشان‌دهنده‌ی یک آرشیو «خام» یا ساده و پسوند دوم معرف آرشیوی است که با ابزار \cmd{gzip} فشرده شده است. یک فایل آرشیو \cmd{tar} می‌تواند شامل مجموعه‌ای از فایل‌های مجزا، یک یا چند ساختار درختی دایرکتوری و یا ترکیبی از هر دو باشد. ساختار کلی این دستور به شرح زیر است:

\begin{latin}
\begin{lstlisting}
tar mode[options] pathname...
\end{lstlisting}
\end{latin}

در این ساختار، پارامتر \file{mode} تعیین‌کننده‌ی نوع عملیات است. جدول زیر برخی از پرکاربردترین حالت‌های عملیاتی \cmd{tar} را معرفی می‌کند (برای مشاهده فهرست کامل، می‌توانید به مستندات سیستم یا \cmd{man tar} مراجعه کنید):

\begin{table}[htbp]
\centering
\caption{حالت‌های عملیاتی (\en{Modes}) دستور \cmd{tar}}
\label{tab:tar-modes}
\begin{tabularx}{\textwidth}{c X}
\toprule
\textbf{حالت (\en{Mode})} & \textbf{شرح عملکرد فنی} \\
\midrule
\cmd{c} & ایجاد یک آرشیو جدید از فایل‌ها و دایرکتوری‌های مشخص‌شده. \\
\addlinespace
\cmd{x} & استخراج محتویات یک آرشیو موجود بر روی دیسک. \\
\addlinespace
\cmd{r} & افزودن فایل‌ها یا دایرکتوری‌های جدید به انتهای یک آرشیو از پیش ایجادشده. \\
\addlinespace
\cmd{t} & فهرست کردن محتویات داخل یک فایل آرشیو بدون استخراج آن‌ها. \\
\bottomrule
\end{tabularx}
\end{table}

نحوه‌ی به‌کارگیری آرگومان‌ها در \cmd{tar} کمی متفاوت از سایر دستورات لینوکس است؛ بنابراین برای درک بهتر، از چند مثال عملی شروع می‌کنیم. ابتدا محیط تمرینی فصل قبل را بازسازی می‌کنیم:

\begin{latin}
\begin{lstlisting}
[me@linuxbox ~]$ mkdir -p playground/dir-{001..100}
[me@linuxbox ~]$ touch playground/dir-{001..100}/file-{A..Z}
\end{lstlisting}
\end{latin}

حالا یک فایل آرشیو از کل دایرکتوری \file{playground} ایجاد می‌کنیم:

\begin{latin}
\begin{lstlisting}
[me@linuxbox ~]$ tar cf playground.tar playground
\end{lstlisting}
\end{latin}

این دستور آرشیوی به نام \file{playground.tar} می‌سازد که شامل تمامی زیرشاخه‌ها و فایل‌های درون آن است. توجه داشته باشید که در اینجا حالت عملیاتی (\cmd{c}) و گزینه‌ی تعیین نام فایل (\cmd{f}) در کنار هم قرار گرفته‌اند و لزوماً نیازی به پیشوند خط تیره (\en{Dash}) ندارند. با این حال، در این سبک نوشتار، حالت عملیاتی همواره باید در ابتدای رشته‌ی گزینه‌ها قرار گیرد.

برای مشاهده‌ی فهرست محتویات آرشیو ساخته شده، از دستور زیر استفاده می‌کنیم:

\begin{latin}
\begin{lstlisting}
[me@linuxbox ~]$ tar tf playground.tar
\end{lstlisting}
\end{latin}

جهت دریافت خروجی با جزئیات بیشتر (مشابه دستور \cmd{ls -l})، می‌توان گزینه \cmd{v} (مخفف \en{verbose}) را به آن افزود:

\begin{latin}
\begin{lstlisting}
[me@linuxbox ~]$ tar tvf playground.tar
\end{lstlisting}
\end{latin}

اکنون برای استخراج آرشیو در یک مسیر جدید، دایرکتوری موسوم به \file{foo} را ایجاد کرده، وارد آن می‌شویم و عملیات استخراج را اجرا می‌کنیم:

\begin{latin}
\begin{lstlisting}
[me@linuxbox ~]$ mkdir foo
[me@linuxbox ~]$ cd foo
[me@linuxbox foo]$ tar xf ../playground.tar
[me@linuxbox foo]$ ls
playground
\end{lstlisting}
\end{latin}

با بررسی محتویات \file{foo/playground/} متوجه خواهید شد که سلسله‌مراتب فایل‌ها دقیقاً مشابه نسخه‌ی اصلی بازیابی شده است.

یک نکته‌ی فنی حائز اهمیت: هنگام استخراج آرشیو توسط کاربر عادی، مالکیت (\en{Ownership}) فایل‌های بازیابی‌شده به کاربری که دستور را اجرا کرده تخصیص می‌یابد؛ تنها در صورتی که با دسترسی ابرکاربر اقدام کنید، مالکیت اصلی فایل‌ها حفظ می‌شود.

رفتار کلیدی دیگر در \cmd{tar}، نحوه‌ی مدیریت نام‌های مسیر (\en{Pathnames}) است. این ابزار به‌طور پیش‌فرض مسیرها را به صورت نسبی (\en{Relative}) ذخیره می‌کند. هنگام ایجاد آرشیو، \cmd{tar} پیشوند اسلش (/) ابتدایی را از مسیرهای مطلق حذف می‌کند تا از بروز خطرات امنیتی و بازنویسی ناخواسته‌ی فایل‌های سیستم در زمان استخراج جلوگیری شود.

برای تحلیل دقیق‌تر این سازوکار، آرشیو را مجدداً با استفاده از مسیر مطلق (\en{Absolute Path}) ایجاد می‌کنیم:

\begin{latin}
\begin{lstlisting}
[me@linuxbox foo]$ cd
[me@linuxbox ~]$ tar cf playground2.tar ~/playground
\end{lstlisting}
\end{latin}

توجه داشته باشید که کاراکتر تیلدا (\cmd{\textasciitilde}) در زمان اجرا توسط پوسته به مسیر کامل \file{/home/me/playground} بسط می‌یابد. حال آرشیو را در دایرکتوری \file{foo} استخراج کرده و ساختار خروجی را بررسی می‌کنیم:

\begin{latin}
\begin{lstlisting}
[me@linuxbox ~]$ cd foo
[me@linuxbox foo]$ tar xf ../playground2.tar
[me@linuxbox foo]$ ls
home    playground
[me@linuxbox foo]$ ls home
me
[me@linuxbox foo]$ ls home/me
playground
\end{lstlisting}
\end{latin}

مشاهده می‌کنید که \cmd{tar} هنگام استخراج، کل سلسله‌مراتب \file{home/me/playground} را به‌صورت نسبی و در دایرکتوری کاری فعلی (یعنی \file{foo/}) بازسازی کرده است، نه در دایرکتوری ریشه (\cmd{/}). این رفتار که با حذف اسلش ابتدایی مسیرها صورت می‌گیرد، انعطاف‌پذیری فوق‌العاده‌ای به مدیران سیستم می‌دهد؛ چرا که امکان بازیابی آرشیوها را در هر مسیر دلخواهی فراهم می‌سازد، بدون آنکه چیدمان اصلی فایل‌های سیستمی را به خطر بیندازد. استفاده از گزینه‌ی \cmd{v} (\en{Verbose}) در زمان ایجاد آرشیو، حذف این اسلش‌ها را به وضوح نمایش می‌دهد.

در ادامه، یک سناریوی عملیاتی را بررسی می‌کنیم: فرض کنید قصد دارید دایرکتوری \file{/home} را از یک سیستم به سیستم دیگر منتقل کنید و برای این کار از یک هارد اکسترنال با برچسب \en{BigDisk} استفاده می‌کنید که در مسیر \file{/media/BigDisk/} متصل (\en{Mount}) شده است. برای ایجاد این آرشیو جامع، دستور زیر را با دسترسی ابرکاربر (\en{superuser}) اجرا می‌کنیم:

\begin{latin}
\begin{lstlisting}
[me@linuxbox ~]$ sudo tar cf /media/BigDisk/home.tar /home
\end{lstlisting}
\end{latin}

در این مثال، استفاده از \cmd{sudo} الزامی است، زیرا دایرکتوری \file{/home} حاوی فایل‌های شخصی تمامی کاربران است که کاربر عادی لینوکس مجاز به خواندن همه‌ی آن‌ها نیست. پس از اتمام، فایل \file{home.tar} حاوی کپی کامل و دقیق از تمامی حساب‌های کاربری سیستم شما خواهد بود که آماده‌ی انتقال به مقصد است.

پس از اتمام فرآیند نوشتن فایل \cmd{tar} بر روی حافظه‌ی جانبی، درایو را از سیستم مبدأ جدا کرده و به کامپیوتر مقصد متصل می‌کنیم. با فرض اینکه درایو مجدداً در مسیر \file{/media/BigDisk/} شناسایی و متصل (\en{Mount}) شده است، عملیات استخراج آرشیو را به شرح زیر اجرا می‌کنیم:

\begin{latin}
\begin{lstlisting}
[me@linuxbox2 ~]$ cd /
[me@linuxbox2 /]$ sudo tar xf /media/BigDisk/home.tar
\end{lstlisting}
\end{latin}

نکته‌ی کلیدی در این مرحله، تغییر دایرکتوری کاری به ریشه (\cmd{/}) پیش از اجرای دستور است. از آنجا که تمام مسیرها در فایل آرشیو به صورت نسبی ذخیره شده‌اند، استخراج باید نسبت به دایرکتوری ریشه صورت گیرد تا فایل‌ها دقیقاً در جایگاه اصلی خود قرار گیرند.

علاوه بر استخراج کامل، ابزار \cmd{tar} امکان بازیابی گزینشی فایل‌ها را نیز فراهم می‌کند. برای استخراج یک یا چند فایل خاص از دل آرشیو، کافی است مسیر دقیق آن‌ها را در انتهای دستور ذکر کنید:

\begin{latin}
\begin{lstlisting}
tar xf archive.tar pathname
\end{lstlisting}
\end{latin}

باید توجه داشت که مسیر اعلام شده باید دقیقاً با مسیر نسبی ثبت‌شده در آرشیو منطبق باشد. در نسخه‌ی استاندارد \cmd{tar}، استفاده از نویسه‌های جانشین (\en{Wildcards}) مجاز نیست؛ اما در نسخه‌ی \en{GNU tar} (که توزیع‌های لینوکس به صورت پیش‌فرض از آن بهره می‌برند)، می‌توان با استفاده از گزینه‌ی \cmd{--wildcards} این قابلیت را فعال کرد. مثال زیر نحوه‌ی استخراج فایل‌های خاص از آرشیو \file{playground2.tar} را نشان می‌دهد:

\begin{latin}
\begin{lstlisting}
[me@linuxbox ~]$ cd foo
[me@linuxbox foo]$ tar xf ../playground2.tar --wildcards 'home/me/playground/dir-*/file-A'
\end{lstlisting}
\end{latin}

در این دستور، تنها فایل‌هایی بازگردانی می‌شوند که با الگوی مسیر \cmd{dir-*/file-A} مطابقت داشته باشند؛ به این ترتیب از استخراج هزاران فایل غیرضروری جلوگیری می‌شود.

ترکیب ابزار \cmd{tar} با دستور \cmd{find} راهکاری بسیار قدرتمند برای ایجاد آرشیوهای هدفمند و مدیریت‌شده است. در مثال زیر، از \cmd{find} جهت شناسایی فایل‌های خاص و درج آن‌ها در یک آرشیو استفاده شده است:

\begin{latin}
\begin{lstlisting}
[me@linuxbox ~]$ find playground -name 'file-A' -exec tar rf playground.tar '{}' '+'
\end{lstlisting}
\end{latin}

در این سناریو، \cmd{find} تمام فایل‌های موسوم به \file{file-A} را در دایرکتوری \file{playground} شناسایی کرده و از طریق عملگر \cmd{-exec} به ابزار \cmd{tar} تحویل می‌دهد. استفاده از حالت \cmd{r} (\en{Append}) باعث می‌شود فایل‌های یافت‌شده به انتهای آرشیو موجود اضافه شوند.

بهره‌گیری از این ترکیب، روشی ایده‌آل برای پیاده‌سازی پشتیبان‌گیری افزایشی (\en{Incremental Backup}) از یک دایرکتوری یا کل سیستم فایل است. با استفاده از شرط \cmd{-newer} در دستور \cmd{find}، می‌توان تنها فایل‌هایی را که پس از آخرین عملیات پشتیبان‌گیری تغییر یافته‌اند شناسایی و آرشیو کرد؛ مشروط بر اینکه پس از هر بار اجرای دستور، فایل مرجع (\en{Timestamp}) به‌روزرسانی شود.

ابزار \cmd{tar} همچنین قابلیت تعامل کامل با ورودی و خروجی استاندارد (\en{Standard I/O}) را داراست. مثال زیر یک زنجیره‌ی عملیاتی جامع را نمایش می‌دهد:

\begin{latin}
\begin{lstlisting}
[me@linuxbox foo]$ cd
[me@linuxbox ~]$ find playground -name 'file-A' | tar cf - --files-from=- | gzip > playground.tgz
\end{lstlisting}
\end{latin}

در این مثال، ابتدا فهرستی از فایل‌های منطبق توسط \cmd{find} تولید و به \cmd{tar} هدایت شده است. مطابق با یک قرارداد مرسوم در ابزارهای خط فرمان لینوکس، نویسه‌ی \cmd{-} به معنای ورودی یا خروجی استاندارد تلقی می‌شود. گزینه‌ی \cmd{--files-from} (یا فرم کوتاه آن \cmd{-T}) به \cmd{tar} فرمان می‌دهد که لیست مسیرها را به جای آرگومان‌های خط فرمان، از یک فایل (که در اینجا همان ورودی استاندارد است) بخواند. در نهایت، خروجی آرشیو به \cmd{gzip} پایپ شده تا فایل فشرده‌ی نهایی با پسوند \file{.tgz} تولید شود. (استفاده از پسوند \file{.tar.gz} نیز برای این نوع فایل‌ها بسیار متداول است).

اگرچه در مثال پیشین از ابزار \cmd{gzip} به‌صورت مجزا برای تولید آرشیو فشرده استفاده شد، نسخه‌های مدرن \en{GNU tar} با بهره‌گیری از گزینه‌های داخلی \cmd{z} و \cmd{j}، به‌طور مستقیم از فشرده‌سازی \cmd{gzip} و \cmd{bzip2} پشتیبانی می‌کنند. بر این اساس، می‌توان مثال قبلی را به شکلی بهینه‌تر بازنویسی کرد:

\begin{latin}
\begin{lstlisting}
[me@linuxbox ~]$ find playground -name 'file-A' | tar czf playground.tgz -T -
\end{lstlisting}
\end{latin}

در صورتی که نیاز به استفاده از الگوریتم \cmd{bzip2} جهت دستیابی به نرخ فشرده‌سازی بالاتر باشد، دستور به شکل زیر تغییر می‌یابد:

\begin{latin}
\begin{lstlisting}
[me@linuxbox ~]$ find playground -name 'file-A' | tar cjf playground.tbz -T -
\end{lstlisting}
\end{latin}

با جایگزینی ساده‌ی پارامتر \cmd{z} با \cmd{j} (و تغییر پسوند فایل خروجی به \file{.tbz} جهت شناسایی نوع فشرده‌سازی)، فرآیند \cmd{bzip2} فعال می‌شود.

یکی از کاربردهای پیشرفته و کارآمد ورودی/خروجی استاندارد در دستور \cmd{tar}، انتقال مستقیم فایل‌ها بین دو سیستم از طریق شبکه است. در صورتی که به دو ماشین با سیستم‌عامل شبه‌لینوکس دسترسی داشته باشید که هر دو به ابزارهای \cmd{tar} و \cmd{ssh} مجهز هستند، می‌توانید یک دایرکتوری را به این صورت منتقل کنید:

\begin{latin}
\begin{lstlisting}
[me@linuxbox ~]$ mkdir remote-stuff
[me@linuxbox ~]$ cd remote-stuff
[me@linuxbox remote-stuff]$ ssh remote-sys 'tar cf - Documents' | tar xf -
me@remote-sys's password:
[me@linuxbox remote-stuff]$ ls
Documents
\end{lstlisting}
\end{latin}

در این سناریو، دایرکتوری \file{Documents} از سیستم راه دور (\file{remote-sys}) به دایرکتوری محلی \file{remote-stuff} کپی شده است. تحلیل فنی این عملیات بدین شرح است: ابتدا دستور \cmd{tar} از طریق پروتکل \cmd{ssh} در سیستم مقصد فراخوانی می‌شود. از آنجا که \cmd{ssh} خروجی استاندارد (\en{stdout}) دستور اجرا شده در سیستم راه دور را به سیستم محلی منتقل می‌کند، ما از \cmd{tar} خواستیم که آرشیو را در حالت \cmd{c} ایجاد کرده و با استفاده از پارامتر \cmd{-f -} مستقیماً به خروجی استاندارد بفرستد. این جریان داده از طریق تونل رمزنگاری‌شده‌ی \cmd{ssh} به سیستم محلی جاری می‌شود و در آنجا توسط دستور \cmd{tar xf -} از ورودی استاندارد (\en{stdin}) دریافت و بلافاصله استخراج می‌گردد.

\section{بسته‌بندی و فشرده‌سازی با ابزار zip}

برنامه‌ی \cmd{zip} ابزاری دو منظوره است که هم‌زمان قابلیت‌های فشرده‌سازی و آرشیو‌سازی را ارائه می‌دهد. قالب فایلی مورد استفاده‌ی این برنامه برای کاربران سیستم‌عامل ویندوز کاملاً شناخته شده است، زیرا ویندوز به‌صورت بومی قادر به خواندن و نوشتن فایل‌های \file{.zip} است. با این وجود، در اکوسیستم لینوکس، ابزار \cmd{gzip} استاندارد غالب محسوب می‌شود و \cmd{bzip2} در جایگاه بعدی قرار دارد.

در ساده‌ترین حالت، نحوه‌ی فراخوانی دستور \cmd{zip} به شرح زیر است:

\begin{latin}
\begin{lstlisting}
zip options zipfile file...
\end{lstlisting}
\end{latin}

برای نمونه، جهت ایجاد یک آرشیو \cmd{zip} از دایرکتوری \file{playground}، دستور زیر را اجرا می‌کنیم:

\begin{latin}
\begin{lstlisting}
[me@linuxbox ~]$ zip -r playground.zip playground
\end{lstlisting}
\end{latin}

نکته‌ی فنی حائز اهمیت این است که برای آرشیو کردن یک دایرکتوری، حتماً باید از پارامتر \cmd{-r} (مخفف \en{Recursive}) استفاده شود؛ در غیر این صورت، تنها خودِ دایرکتوری \file{playground} (به‌صورت یک ورودی خالی) ذخیره شده و محتویات درون آن نادیده گرفته می‌شود. اگرچه درج پسوند \file{.zip} توسط برنامه به‌صورت خودکار انجام می‌گردد، اما برای شفافیت بیشتر، آن را در دستور لحاظ می‌کنیم.

در طول فرآیند ایجاد آرشیو، ابزار \cmd{zip} گزارش‌هایی مشابه نمونه‌ی زیر را در خروجی نمایش می‌دهد:

\begin{latin}
\begin{lstlisting}
adding: playground/dir-020/file-Z (stored 0%)
adding: playground/dir-020/file-Y (stored 0%)
adding: playground/dir-020/file-X (stored 0%)
adding: playground/dir-087/ (stored 0%)
adding: playground/dir-087/file-S (stored 0%)
\end{lstlisting}
\end{latin}

این پیام‌ها بیانگر وضعیت هر فایل در هنگام الحاق به آرشیو هستند. ابزار \cmd{zip} به دو روش فایل‌ها را مدیریت می‌کند: یا فایل را بدون فشرده‌سازی صرفاً «ذخیره» (\en{Stored}) می‌کند و یا با استفاده از الگوریتم فشرده‌سازی، آن را فشرده (\en{Deflated}) می‌کند. عدد درصدی که پس از روش ذخیره‌سازی درج می‌شود، نشان‌دهنده‌ی میزان کاهش حجم فایل است. در مثال فوق، از آنجا که دایرکتوری \file{playground} تنها شامل فایل‌های تهی است، هیچ فشرده‌سازی مؤثری صورت نگرفته و مقدار آن صفر درصد درج شده است.

استخراج محتوای یک فایل \cmd{zip} با استفاده از ابزار \cmd{unzip} به سهولت امکان‌پذیر است:

\begin{latin}
\begin{lstlisting}
[me@linuxbox ~]$ cd foo
[me@linuxbox foo]$ unzip ../playground.zip
\end{lstlisting}
\end{latin}

یک تفاوت ساختاری مهم میان \cmd{zip} و \cmd{tar} در نحوه‌ی مواجهه با آرشیوهای موجود است؛ ابزار \cmd{zip} به جای جایگزینی (\en{Overwrite}) کامل فایل آرشیو، آن را به‌روزرسانی (\en{Update}) می‌کند. این بدان معناست که فایل قبلی حفظ شده، فایل‌های جدید به آن ضمیمه می‌گردند و تنها فایل‌های هم‌نام جایگزین می‌شوند.

همچنین، می‌توان فایل‌ها را به‌صورت انتخابی از یک آرشیو \cmd{zip} فهرست کرده یا استخراج نمود:

\begin{latin}
\begin{lstlisting}
[me@linuxbox ~]$ unzip -l playground.zip playground/dir-087/file-Z
Archive:  ../playground.zip
  Length     Date   Time    Name
---------    ----   ----    ----
        0  10-05-16 09:25   playground/dir-087/file-Z
---------                   -------
        0                   1 file

[me@linuxbox ~]$ cd foo
[me@linuxbox foo]$ unzip ../playground.zip playground/dir-087/file-Z
Archive:  ../playground.zip
replace playground/dir-087/file-Z? [y]es, [n]o, [A]ll, [N]one, [r]ename: y
extracting: playground/dir-087/file-Z
\end{lstlisting}
\end{latin}

استفاده از پارامتر \cmd{-l} به ابزار \cmd{unzip} فرمان می‌دهد که صرفاً محتویات آرشیو را فهرست‌بندی کند، بدون آنکه عملیات استخراج را انجام دهد. در صورتی که نام فایل خاصی قید نشود، تمامی فایل‌های موجود در آرشیو لیست خواهند شد. همچنین با افزودن گزینه \cmd{-v}، می‌توان جزئیات فنی بیشتری را در این فهرست مشاهده کرد. شایان ذکر است که اگر هنگام استخراج، فایلی هم‌نام در مقصد وجود داشته باشد، \cmd{unzip} پیش از بازنویسی (\en{Overwrite})، از کاربر تاییدیه می‌گیرد.

مشابه ابزار \cmd{tar}، برنامه‌ی \cmd{zip} نیز قابلیت تعامل با جریان‌های استاندارد ورودی و خروجی (\en{standard I/O streams}) را دارد، هرچند پیاده‌سازی آن محدودتر است. برای نمونه، می‌توان فهرستی از نام فایل‌ها را از طریق پارامتر \cmd{-@} به ورودی \cmd{zip} هدایت کرد:

\begin{latin}
\begin{lstlisting}
[me@linuxbox foo]$ cd
[me@linuxbox ~]$ find playground -name "file-A" | zip -@ file-A.zip
\end{lstlisting}
\end{latin}

در این مثال، خروجی دستور \cmd{find} به عنوان لیست ورودی به \cmd{zip} تحویل داده شده و آرشیو \file{file-A.zip} تنها شامل فایل‌های منطبق ایجاد می‌گردد.

اگرچه \cmd{zip} توانایی ارسال خروجی به \en{stdout} را دارد، اما به دلیل محدودیت‌های فنی و عدم پشتیبانی \cmd{unzip} از ورودی استاندارد (\en{stdin})، نمی‌توان از ترکیب این دو برای کپی مستقیم فایل‌ها در شبکه (مشابه روش \cmd{tar | ssh}) بهره برد. با این حال، \cmd{zip} می‌تواند خروجی سایر برنامه‌ها را فشرده‌سازی کند:

\begin{latin}
\begin{lstlisting}
[me@linuxbox ~]$ ls -l /etc/ | zip ls-etc.zip -
 adding: - (deflated 80%)
\end{lstlisting}
\end{latin}

در اینجا، نویسه‌ی خط تیره (\cmd{-}) در انتهای دستور به معنای پذیرش ورودی از ورودی استاندارد (\en{stdin}) است. متقابلاً، ابزار \cmd{unzip} با استفاده از پارامتر \cmd{-p} (مخفف \en{pipe})، محتوای فشرده را مستقیماً به خروجی استاندارد ارسال می‌کند:

\begin{latin}
\begin{lstlisting}
[me@linuxbox ~]$ unzip -p ls-etc.zip | less
\end{lstlisting}
\end{latin}

ما در اینجا تنها به قابلیت‌های کلیدی و پایه پرداخته‌ایم. هر دو برنامه‌ی \cmd{zip} و \cmd{unzip} دارای پارامترهای عملیاتی متعددی هستند که انعطاف‌پذیری آن‌ها را ارتقا می‌دهد؛ هرچند برخی از این گزینه‌ها صرفاً برای سازگاری با سیستم‌عامل‌های غیرلینوکسی تعبیه شده‌اند.

در نهایت باید توجه داشت که کاربرد اصلی این دو برنامه در دنیای لینوکس، تسهیل تبادل فایل با محیط‌های ویندوزی است؛ چرا که برای مقاصد بومی لینوکس، ابزارهای \cmd{tar}، \cmd{gzip} و \cmd{bzip2} به دلیل یکپارچگی بیشتر با معماری سیستم، همواره در اولویت قرار دارند.

\section{همگام‌سازی فایل‌ها و دایرکتوری‌ها}

یکی از راهبردهای استاندارد برای مدیریت داده‌ها و نگهداری نسخه‌های پشتیبان، همگام‌سازی (\en{Synchronization}) یک یا چند دایرکتوری با مجموعه‌ای دیگر از دایرکتوری‌ها است. این دایرکتوری‌های مقصد می‌توانند در همان سیستم محلی (مانند یک حافظه‌ی جانبی) و یا در یک سرور راه دور مستقر باشند. برای مثال، یک توسعه‌دهنده ممکن است نسخه‌ی محلی وب‌سایت خود را به‌طور منظم با نسخه‌ی عملیاتی (\en{Live}) موجود بر روی یک سرور راه دور همگام‌سازی کند.

در اکوسیستم سیستم‌عامل‌های شبه‌یونیکس، ابزار استاندارد و مرجع برای این عملیات، \cmd{rsync} است. این برنامه قادر است دایرکتوری‌های محلی و راه دور را از طریق «الگوریتم به‌روزرسانی راه دور \cmd{rsync}» همگام‌سازی نماید. ویژگی متمایز این الگوریتم، توانایی شناسایی سریع تفاوت‌های میان دو دایرکتوری و انتقال کمترین حجم ممکن از داده‌ها برای یکسان‌سازی آن‌هاست. این سازوکار باعث شده تا \cmd{rsync} در مقایسه با سایر ابزارهای کپی‌برداری، بسیار سریع‌تر و از نظر مصرف پهنای باند و منابع، بهینه‌تر عمل کند.

ساختار کلی فراخوانی دستور \cmd{rsync} به شرح زیر است:

\begin{latin}
\begin{lstlisting}
rsync options source destination
\end{lstlisting}
\end{latin}

در این ساختار، مبدأ (\en{source}) و مقصد (\en{destination}) می‌توانند شامل موارد زیر باشند:

\begin{itemize}
  \item یک فایل یا دایرکتوری در سیستم محلی.
  \item یک فایل یا دایرکتوری در سیستم راه دور که به فرمت زیر مشخص می‌شود:
\begin{latin}
\begin{lstlisting}
[user@]host:path
\end{lstlisting}
\end{latin}
  \item یک سرور راه دور \cmd{rsync} که از طریق \en{URI} و به فرمت ذیل تعریف می‌گردد:
\begin{latin}
\begin{lstlisting}
rsync://[user@]host[:port]/path
\end{lstlisting}
\end{latin}
\end{itemize}

نکته‌ی فنی حائز اهمیت این است که در هر دستور، دست‌کم یکی از طرفین (مبدأ یا مقصد) باید یک مسیر محلی باشد؛ چرا که \cmd{rsync} از کپی‌برداری مستقیم بین دو سرور راه دور (\en{Remote-to-Remote}) پشتیبانی نمی‌کند.

بیایید عملکرد \cmd{rsync} را بر روی فایل‌های محلی بررسی کنیم. ابتدا دایرکتوری \file{foo} را برای شروع یک فرآیند تازه پاکسازی می‌کنیم:

\begin{latin}
\begin{lstlisting}
[me@linuxbox ~]$ rm -rf foo/*
\end{lstlisting}
\end{latin}

سپس دایرکتوری \file{playground} را در مسیر \file{foo} همگام‌سازی (\en{Mirror}) می‌کنیم:

\begin{latin}
\begin{lstlisting}
[me@linuxbox ~]$ rsync -av playground foo
\end{lstlisting}
\end{latin}

در این دستور، از پارامتر \cmd{-a} (مخفف \en{Archive}) برای پیمایش بازگشتی دایرکتوری‌ها و حفظ تمامی ویژگی‌های فایل (مانند مجوزها، مالکیت و زمان‌ها) و از گزینه \cmd{-v} (مخفف \en{Verbose}) جهت مشاهده‌ی جزئیات عملیات استفاده کردیم.

پس از اجرا، فهرستی از فایل‌های در حال انتقال نمایش داده می‌شود و در پایان، گزارشی آماری ارائه می‌گردد که بیانگر حجم داده‌های تبادل شده است:

\begin{latin}
\begin{lstlisting}
sent 135759 bytes  received 57870 bytes  387258.00 bytes/sec
total size is 3230  speedup is 0.02
\end{lstlisting}
\end{latin}

چنانچه بلافاصله دستور را مجدداً اجرا کنیم، خروجی متفاوتی دریافت خواهیم کرد:

\begin{latin}
\begin{lstlisting}
[me@linuxbox ~]$ rsync -av playground foo
building file list ... done

sent 22635 bytes  received 20 bytes  45310.00 bytes/sec
total size is 3230  speedup is 0.14
\end{lstlisting}
\end{latin}

مشاهده می‌کنید که این بار هیچ فایلی کپی نشد؛ چرا که \cmd{rsync} با بررسی هوشمند متوجه شد که محتویات مبدأ و مقصد کاملاً یکسان هستند. حال اگر تغییری در یکی از فایل‌های مبدأ ایجاد کنیم:

\begin{latin}
\begin{lstlisting}
[me@linuxbox ~]$ touch playground/dir-099/file-Z
[me@linuxbox ~]$ rsync -av playground foo
building file list ... done
playground/dir-099/file-Z
sent 22685 bytes  received 42 bytes  45454.00 bytes/sec
total size is 3230  speedup is 0.14
\end{lstlisting}
\end{latin}

در این مرحله، \cmd{rsync} تنها فایل تغییر یافته را شناسایی و منتقل کرد که نشان‌دهنده‌ی کارایی بالای این ابزار در مدیریت داده‌هاست.

یک نکته‌ی ظریف اما بسیار کاربردی در تعیین مسیر مبدأ در \cmd{rsync} وجود دارد که رفتار این ابزار را تغییر می‌دهد. دو دایرکتوری فرضی زیر را در نظر بگیرید:

\begin{latin}
\begin{lstlisting}
[me@linuxbox ~]$ ls
source      destination
\end{lstlisting}
\end{latin}

فرض کنید دایرکتوری \file{source} حاوی فایلی به نام \file{file1} است و \file{destination} خالی است. اگر دستور را به شکل زیر اجرا کنید:

\begin{latin}
\begin{lstlisting}
[me@linuxbox ~]$ rsync source destination
\end{lstlisting}
\end{latin}

در این حالت، \cmd{rsync} خودِ دایرکتوری \file{source} را به درون مقصد کپی می‌کند:

\begin{latin}
\begin{lstlisting}
[me@linuxbox ~]$ ls destination
source
\end{lstlisting}
\end{latin}

اما اگر یک اسلش (\cmd{/}) به انتهای نام دایرکتوری مبدأ اضافه کنید، \cmd{rsync} تنها محتویات دایرکتوری مبدأ را کپی کرده و از ایجاد دایرکتوری والد در مقصد صرف‌نظر می‌کند:

\begin{latin}
\begin{lstlisting}
[me@linuxbox ~]$ rsync source/ destination
[me@linuxbox ~]$ ls destination
file1
\end{lstlisting}
\end{latin}

این روش زمانی بسیار کارآمد است که قصد دارید محتویات یک دایرکتوری را بدون ایجاد لایه‌های اضافی در ساختار درختی مقصد منتقل کنید. این رفتار مشابه استفاده از نویسه‌ی جانشین (\cmd{source/*}) است، با این تفاوت که در روش اسلش انتهایی، تمامی فایل‌ها از جمله فایل‌های پنهان (\en{Dotfiles}) نیز به دقت کپی می‌شوند.

به‌عنوان یک سناریوی عملی، همان هارد اکسترنال فرضی را در نظر بگیرید که در مسیر \file{/media/BigDisk/} متصل (\en{Mount}) شده است. می‌توان با ایجاد یک دایرکتوری پشتیبان و استفاده از \cmd{rsync}، یک کپی از بخش‌های حیاتی سیستم تهیه کرد:

\begin{latin}
\begin{lstlisting}
[me@linuxbox ~]$ mkdir /media/BigDisk/backup
[me@linuxbox ~]$ sudo rsync -av --delete /etc /home /usr/local /media/BigDisk/backup
\end{lstlisting}
\end{latin}

در این فرمان، دایرکتوری‌های \file{/etc}، \file{/home} و \file{/usr/local} به حافظه‌ی جانبی منتقل می‌شوند. استفاده از پارامتر \cmd{--delete} تضمین می‌کند که اگر فایلی در سیستم اصلی حذف شده باشد، در نسخه‌ی پشتیبان نیز حذف گردد تا دو محیط کاملاً همگام باقی بمانند (این گزینه در اجرای اول تأثیری ندارد اما در دفعات بعدی برای حفظ یکپارچگی حیاتی است).

تکرار دوره‌ای این فرآیند، راهکاری ساده و مؤثر برای پشتیبان‌گیری از سیستم‌های کوچک و ایستگاه‌های کاری (\en{Workstations}) محسوب می‌شود.

بدون شک، تعریف یک «نام مستعار» (\en{alias}) در این سناریو بسیار کارآمد خواهد بود. برای تسهیل در اجرای فرآیند، می‌توان یک \cmd{alias} اختصاصی ایجاد و آن را در فایل پیکربندی \file{.bashrc} درج نمود تا این قابلیت همواره در دسترس باشد:

\begin{latin}
\begin{lstlisting}
alias backup='sudo rsync -av --delete /etc /home /usr/local /media/BigDisk/backup'
\end{lstlisting}
\end{latin}

با اعمال این تنظیم، از این پس تنها کافی است پس از اتصال هارددیسک خارجی (\en{External Drive})، فرمان کوتاه \cmd{backup} را در ترمینال اجرا کنید تا عملیات همگام‌سازی و پشتیبان‌گیری به‌صورت خودکار آغاز گردد.

\subsection{همگام‌سازی تحت شبکه با ابزار rsync}

یکی از قابلیت‌های برجسته و کلیدی \cmd{rsync} امکان انتقال و همگام‌سازی فایل‌ها در بستر شبکه است؛ در واقع پیشوند «\en{r}» در نام این ابزار، به واژه‌ی «\en{remote}» اشاره دارد. عملیات انتقال داده در شبکه به دو شیوه‌ی مجزا قابل پیاده‌سازی است.

روش نخست، مبتنی بر سیستمی است که علاوه بر نصب بودن \cmd{rsync} بر روی آن، از یک پوسته راه دور (\en{Remote Shell}) مانند \cmd{ssh} نیز بهره می‌برد. تصور کنید در شبکه‌ی محلی خود، سیستمی با فضای ذخیره‌سازی انبوه در اختیار دارید و قصد دارید به جای درایو خارجی، از آن به عنوان مقصد پشتیبان‌گیری استفاده کنید. اگر در آن سیستم راه دور، دایرکتوری \file{/backup} برای میزبانی فایل‌های شما آماده باشد، دستور زیر عملیات انتقال را انجام می‌دهد:

\begin{latin}
\begin{lstlisting}
[me@linuxbox ~]$ sudo rsync -av --delete --rsh=ssh /etc /home /usr/local remote-sys:/backup
\end{lstlisting}
\end{latin}

در این ساختار، دو تغییر اساسی اعمال شده است: ابتدا گزینه‌ی \cmd{--rsh=ssh} اضافه شده که به \cmd{rsync} فرمان می‌دهد از پروتکل \cmd{ssh} به عنوان واسط ارتباطی استفاده کند. بدین ترتیب، انتقال داده‌ها از طریق یک تونل رمزنگاری‌شده و ایمن صورت می‌گیرد. دوم، مسیر مقصد با پیشوند نام میزبان راه دور (در اینجا \file{remote-sys}) مشخص شده است تا نقطه پایانی انتقال برای سیستم تعریف شود.

روش دوم برای همگام‌سازی فایل‌ها در شبکه، استفاده از «سرور \cmd{rsync}» است. در این پیکربندی، \cmd{rsync} به صورت یک سرویس پس‌زمینه (\en{Daemon}) اجرا شده و به درخواست‌های ورودی جهت همگام‌سازی گوش فرا می‌دهد. این رویکرد به‌طور معمول برای ایجاد میرورهای محلی (\en{Mirroring}) از مخازن نرم‌افزاری یا سیستم‌های توزیع فایل راه دور در مقیاس وسیع مورد استفاده قرار می‌گیرد.

به‌عنوان نمونه، شرکت \en{Red Hat} مخزن عظیمی از بسته‌های نرم‌افزاری در حال توسعه را برای توزیع \en{Fedora} نگهداری می‌کند. برای آزمون‌گران (آزمایش‌کننده) نرم‌افزار (\en{Testers}) بسیار کارآمدتر است که در طول چرخه‌ی انتشار این توزیع، به جای کپی‌برداری کامل و مکرر از مخزن، یک «نسخه‌ی میرور» (\en{Mirror}) محلی ایجاد کرده و آن را همگام‌سازی کنند. با توجه به اینکه محتوای این مخازن به‌طور مداوم (اغلب چندین بار در روز) دستخوش تغییر می‌شود، نگهداری یک میرور محلی از طریق همگام‌سازی دوره‌ای، رویکردی بسیار بهینه‌تر نسبت به دانلود مجدد کل مخزن است.

یکی از این مخازن در دانشگاه \en{Duke} میزبانی می‌شود؛ ما می‌توانیم با استفاده از کلاینت محلی \cmd{rsync} و اتصال به سرور \cmd{rsync} این دانشگاه، فرآیند میرور کردن را به شکل زیر اجرا کنیم:

\begin{latin}
\begin{lstlisting}
[me@linuxbox ~]$ mkdir fedora-devel
[me@linuxbox ~]$ rsync -av --delete rsync://archive.linux.duke.edu/fedora/linux/development/rawhide/Everything/x86_64/os/ fedora-devel
\end{lstlisting}
\end{latin}

در این دستور، از یک شناسه یکنواخت منبع (\en{URI}) برای سرور راه دور استفاده شده است که با پروتکل اختصاصی \cmd{rsync://} آغاز می‌گردد. متعاقب آن، نام میزبان راه دور (\file{archive.linux.duke.edu}) و در نهایت مسیر دقیق دسترسی به مخزن درج شده است تا عملیات همگام‌سازی با دقت انجام پذیرد.

\section{جمع‌بندی}

در این فصل، ابزارهای کلیدی و پرکاربرد فشرده‌سازی و آرشیو‌سازی (\en{Archiving}) در محیط لینوکس و سایر سیستم‌های شبه‌یونیکس را مورد واکاوی قرار دادیم. آموختیم که برای مدیریت بایگانی‌ها در اکوسیستم‌های یونیکسی، ترکیب \cmd{tar/gzip} استاندارد مرجع و راهکار بهینه محسوب می‌شود؛ در حالی که ترکیب \cmd{zip/unzip} عمدتاً جهت حفظ سازگاری و تبادل داده با محیط‌های ویندوزی به کار می‌رود.

در نهایت، ابزار قدرتمند \cmd{rsync} را بررسی کردیم که به جرئت می‌توان آن را منعطف‌ترین ابزار برای همگام‌سازی هوشمند فایل‌ها و دایرکتوری‌ها بین مقاصد محلی و راه دور دانست.

\section{منابع مطالعه تکمیلی}

صفحات راهنما (\en{man pages}) برای تمامی دستورات مورد بحث، کاملاً واضح بوده و حاوی مثال‌های کاربردی هستند. علاوه بر این، پروژه گنو مستندات آنلاین ارزشمندی برای نسخه خود از ابزار \cmd{tar} ارائه می‌دهد که در آدرس زیر قابل دسترسی است:

\begin{latin}
\url{http://www.gnu.org/software/tar/manual/index.html}
\end{latin}